\chapter{Good and Bad Swing Practice}
\label{ch:practices}
\chapterquote{Swing lets you write anything in a listener; maintainable Swing depends on deciding what must not live there.}{A maintenance warning}

\begin{chaptermap}
Core question & Which common Swing habits fail as screens grow, and what should replace them?\\
Bad smells & Listener sprawl, hidden state, blocking the EDT, view-index persistence, renderer allocation, duplicated resources, focus hacks, and forgotten cleanup.\\
Corusco connection & Corusco names repeated relationships: values, commands, descriptors, problems, bindings, behaviors, scopes, and generated companions.\\
Companion example & \practiceExample.\\
\end{chaptermap}

\section{Practice is architecture in small pieces}

Bad Swing code rarely announces itself as bad architecture. It arrives as one convenient listener, one hard-coded column index, one quick \api{setPreferredSize}, one background thread that updates a label, one tooltip overwrite, one renderer allocation, and one forgotten timer. Each local decision seems cheap. The accumulated screen becomes expensive.

Good practice is not a collection of taboos. It is a collection of ownership rules. Who owns this state? Who owns this listener? Who owns this command? Who owns this validation problem? Who owns the cleanup? If the answer is clear, Swing code tends to remain understandable. If the answer is hidden in component callbacks, the screen decays.

\begin{principlebox}{Review the ownership, not only the syntax}
Modern Java syntax can make poor Swing code shorter. It cannot make ownership clear by itself. A lambda that hides state in a listener is still hidden state.
\end{principlebox}

The best way to read this chapter is as a review method. Each bad practice is a
local symptom of a missing owner. Listener sprawl means no one owns the
relationship graph. Blocking the EDT means no one owns latency. Index-based
table code means no one owns identity. Renderer allocation means no one owns
paint cost. Focus hacks mean no one owns keyboard workflow. Duplicated strings
mean no one owns resources. Leaks mean no one owns lifetime. Once the owner is
named, the remedy usually becomes much less mysterious.

This is also why the chapter belongs early in the book. A reader who learns
only APIs may write code that compiles and still becomes unmaintainable. A
reader who learns ownership can use plain Swing and Corusco selectively. The
goal is not to make every small utility look like an enterprise framework. The
goal is to recognize the point where a local shortcut stops being local.

Every production Swing screen eventually acquires history. A button added for a
demo becomes part of a workflow. A listener written for one field becomes copied
to five fields. A table that had five rows receives ten thousand. A tooltip that
held help text begins to hold validation. A dialog that was modal and simple
gains Apply, dirty confirmation, and background validation. Good practice is the
habit of noticing when the original shape is no longer true.

Corusco should be introduced at those pressure points. If the screen has one
button and one local action, plain Swing may be ideal. If the same operation
appears in a toolbar, menu, shortcut, context menu, and generated test, it wants
a command. If one label displays one local message, plain Swing may be enough.
If validation, status, disabled reasons, and help all compete for the same
tooltip, the screen wants a tooltip policy and problem model. Use the library
where it names real pressure, not where it merely adds vocabulary.

A useful practice review has five passes. The first pass asks what the screen
promises to the user: which operations exist, which data is visible, which
state can change, which errors can block progress, and what happens if the
screen is closed. The second pass asks where those promises live in code. The
third pass asks which Swing mechanics present them: listeners, models,
renderers, actions, input maps, layout constraints, and repaint requests. The
fourth pass asks what the slow path is: data loading, validation, image
decoding, report generation, database access, remote calls, and repaint hot
spots. The fifth pass asks how the relationships end. This order keeps review
from degenerating into a search for personal style preferences.

The passes are deliberately ordinary. They can be applied to a ten-line demo or
to a large generated screen. A small demo may answer each question with a local
object and remain plain Swing. A production screen will often answer with
values, commands, problem sets, descriptors, task handles, and scopes. The
important point is that the same question is asked in both cases. Corusco should
not change the nature of the question; it should make the answer shorter, more
typed, and easier to test when the answer has grown beyond a local component.

This chapter therefore treats every bad habit as a failed boundary. Listener
sprawl is a failed state boundary. Blocking the EDT is a failed latency
boundary. View-index persistence is a failed identity boundary. Renderer
allocation is a failed hot-path boundary. Focus hacks are failed workflow
boundaries. Duplicated resources are failed vocabulary boundaries. Leaks are
failed lifetime boundaries. The naming may sound severe, but it is useful
because it makes remedies comparable. A boundary can be moved, named, tested,
and shared.

\textbf{Listener sprawl.}
Anonymous listeners are not bad. Unowned listeners are bad. A listener that updates one nearby label may be clear. Ten listeners that all partly maintain command state, validation text, table selection, dirty flags, and status messages become a hidden program.

The first warning sign is duplication. Several listeners call the same helper, but each with slightly different assumptions. The second warning sign is ordering. One listener updates a model, another listener observes it, a third listener changes a button, and a fourth listener rewrites a tooltip. The third warning sign is cleanup. No one can say which object removes the listeners.

Listener sprawl is especially deceptive because the first listener is usually
reasonable. A document listener updates a preview. A selection listener enables
Edit. A button listener calls Save. The problem begins when each listener also
knows one extra fact. The document listener updates validation and Save. The
selection listener updates a status label and a detail panel. The button
listener starts a background task and changes a tooltip. Soon there is no single
place to answer "what is the state of this screen?"

A useful review exercise is to draw arrows. Put components on the left, models
and commands on the right, and listeners in the middle. If arrows cross
everywhere, the screen is probably using listeners as its architecture. If most
listeners translate component events into model changes and most component
updates come from model bindings, the screen has a shape a reviewer can follow.
The picture need not be formal; even a rough sketch often reveals the hidden
program.

Do not confuse listener count with listener sprawl. A complex screen may have
many listeners and still be disciplined if they are generated, scoped, and
single-purpose. A tiny screen may have only three listeners and still be
unmaintainable if each one updates several unrelated facts. The question is
whether the listener owns meaning that should live elsewhere.

\begin{migrationbox}{Sprawl to Scope}
Group installed listeners and subscriptions under a lifecycle object. In Corusco Swing, a \api{BindingScope} or \api{BehaviorScope} says, "these relationships belong to this screen activation, and they close together."
\end{migrationbox}

\textbf{Listeners that own facts.}
A listener should translate an event into a model change or command invocation. It should rarely be the sole owner of the fact. A document listener can say "the raw text changed". It should not be the only representation of validity, dirty state, command enablement, status text, and problem routing.

\begin{lstlisting}[caption={A listener that has accumulated too many jobs.}]
nameField.getDocument().addDocumentListener(new SimpleDocumentListener(() -> {
    String text = nameField.getText();
    boolean valid = !text.isBlank();
    errorLabel.setText(valid ? "" : "Name is required");
    saveButton.setEnabled(valid && !busy);
    statusLabel.setText(valid ? "Ready" : "Fix validation errors");
}));
\end{lstlisting}

The listener is readable in isolation, but it owns too many facts. If a menu item also invokes Save, the button's enabled state is the wrong owner. If a dialog summary needs validation, the error label is the wrong owner. If a test needs to assert dirty state, the text field is the wrong owner.

The deeper problem is that the listener has no vocabulary for intermediate
states. It treats text, validity, busy state, command enablement, and status as
nearby widget properties. A real screen needs these facts independently. The
field text may be invalid but touched. The form may be valid but pristine. Save
may be disabled because the form is busy, not because the text is invalid. The
status line may need to show a background failure even while the field is valid.
Once these distinctions exist, a listener is too small a home.

\textbf{The model-backed alternative.}
The alternative is not to eliminate listeners. The alternative is to make them narrow and owned. A binding listens to the document and updates a field model. A validation model exposes problems. A command observes committability. A status value describes the current user-facing message.

The model-backed alternative has a rhythm. First, input changes one model fact:
raw text, selected row, checked state, chosen option. Second, derived state is
computed from named facts: committable, dirty, busy, primary problem, selected
customer name. Third, components display those facts through bindings,
renderers, actions, or behaviors. Fourth, a lifecycle scope closes the
relationships. The rhythm is not elaborate; it is simply explicit.

Tests become shorter because the model facts can be exercised directly. A test
does not need to type into a field to prove that a blank name blocks Save; it
can set the field model and assert the command state. An EDT binding test still
proves that the field edits the model, but it is no longer the only test of the
business rule. This separation is one of the strongest practical reasons to
move facts out of listeners.

The same separation improves debugging. When a user reports that Save is
disabled, the developer should be able to inspect a small chain of facts:
permission, dirty state, parse state, validation problems, busy state, and any
active editor commit failure. If those facts are scattered across buttons,
labels, and listeners, debugging becomes archaeology. If the facts are model
values and command state, a log line or debugger watch can show why the command
is disabled without repainting the whole screen in one's head.

There is also a timing benefit. Component callbacks occur at moments chosen by
Swing: document mutation, focus transfer, selection change, key binding,
timer tick, repaint, property change. Business facts often have their own
meaningful moments: the form became dirty, the form became committable, the
selected customer changed, a background result became stale, a validation
problem was cleared. Conflating the callback moment with the business moment is
one reason listener code grows strange. A model-backed design lets Swing events
enter the model at narrow points, after which the screen can reason in its own
domain terms.

Corusco's observable values and command descriptors are useful precisely
because they make such chains explicit. A \api{WritableValue} can receive
edited text. A derived value can express committability. A command can depend
on that derived value and on busy state. A binding can project command state to
a button, menu item, and shortcut. The old listener has not vanished; it has
been reduced to an adapter from a Swing event to a named fact. The resulting
code is not merely easier to read. It is easier to review because each
relationship has a type and an owner.

When the screen is still small, the same idea can be applied without Corusco.
Create a tiny screen-presenter object with fields for state and methods for events.
Let listeners call those methods. Let the presenter update the few components it
owns. This is already better than hiding the whole screen in anonymous
callbacks. Later, if the screen grows, the presenter methods reveal exactly
which facts want values, commands, validation models, or generated companions.
Good plain Swing can therefore be a migration path rather than a dead end.

\begin{figure}[htbp]
\centering
\includegraphics[width=0.86\linewidth]{\practiceFigure}
\caption{The practice comparison companion example. It is visually small because the point is not widget variety; it is the distinction between component text, model-backed value, and derived status.}
\label{fig:practice-comparison}
\end{figure}

\begin{figure}[htbp]
\centering
\includegraphics[width=0.94\linewidth]{\practicePitfallFigure}
\caption{The review move behind many practice improvements: replace local listener piles with named state, derived command facts, and lifecycle-owned relationships.}
\label{fig:practice-listener-sprawl}
\end{figure}

\begin{examplebox}{Practices}
\practiceExample{} contrasts a local-listener form with a model-backed form and records the review questions each style creates.
\end{examplebox}

\textbf{Blocking the EDT.}
The EDT can perform small calculations. It cannot perform unpredictable I/O, long database queries, slow network calls, or bulk parsing while the user waits for the window to repaint. Modern Java makes background work cheaper; it does not make Swing thread-safe.

The bad version often looks harmless:

\begin{lstlisting}[caption={The freeze is hidden inside the listener.}]
refreshButton.addActionListener(event -> {
    statusLabel.setText("Loading...");
    List<CustomerRow> rows = customerService.loadCustomers(); // Blocks.
    tableModel.setRows(rows);
    statusLabel.setText("Loaded " + rows.size() + " rows");
});
\end{lstlisting}

The label may not even repaint to "Loading..." before the service call blocks the EDT. The user sees a frozen window and perhaps assumes the application has crashed. The correct shape moves unpredictable work out of the EDT and publishes a small result back.

The failure is not merely that the code is slow. It is that the user receives a
lie. The screen says "Loading..." in source code, but the message may never
paint before the window freezes. The button appears clickable but cannot
respond. The window may not repaint when moved. A progress bar cannot progress
because the event queue is occupied by the operation that was supposed to start
progress. Blocking the EDT breaks the conversation between the application and
the user.

Good EDT practice begins by classifying work by latency. Bounded presentation
work belongs on the EDT: setting labels, updating component properties,
installing bindings, replacing a small local model, formatting a short message.
Unbounded or external work does not: file I/O, database queries, remote calls,
large parsing, image decoding, report generation, and loops over user-sized
data. When in doubt, ask whether the operation can become slow because of data,
network, disk, service load, or environment. If yes, keep it out of the EDT.

Modern Java virtual threads reduce the cost of moving blocking work away from
the EDT, but they do not change Swing ownership. The worker may be virtual; the
component mutation still belongs on the EDT. The worker may be cheap to start;
the result may still arrive after the screen closes. The command may submit work
easily; the UI must still model busy state, cancellation, failure, and stale
results.

Latency should be reviewed as part of the screen contract, not as an afterthought
after the first freeze. A button that can start a network operation needs a
busy state before it needs a progress bar. A table that can load a large result
set needs a partial or empty state before it needs a custom renderer. A dialog
that validates against a server needs an answer to the question "what happens
if the user edits again before the old answer returns?" These are practice
questions because they decide where state belongs. The EDT chapter explains the
mechanics; this chapter asks whether the screen has made a promise it can keep.

A small latency budget is useful even when no stopwatch is used. Work that is
effectively immediate may remain in the event handler. Work that may take a
perceptible moment should update visible state first and then leave the EDT.
Work that may outlive the user's patience needs cancellation, progress, or a
clear stale-result rule. Work that may outlive the screen needs lifecycle
ownership. These categories prevent the common mistake of treating all
background work as one topic. A cached calculation, a database query, a remote
upload, and a month-end report do not need the same UI policy.

The visible state matters as much as the worker. Disable or retitle commands
that cannot safely be invoked twice. Keep cancelable tasks cancelable from the
keyboard as well as the mouse. Publish progress at a rate the EDT can consume;
do not replace one freeze with a storm of progress events. Decide whether old
data remains visible while new data loads, whether it is marked stale, or
whether the table moves to an empty loading state. Users can tolerate waiting
when the screen remains truthful. They cannot reason about a window that lies
about being ready.

\begin{pitfallbox}{The Fast Database Fallacy}
The developer says the query is fast. The network says otherwise. The production database says otherwise. The user's VPN says otherwise. EDT policy should be based on worst plausible latency, not the developer's local happy path.
\end{pitfallbox}

\textbf{Unsafe background updates.}
The opposite mistake is also common. A developer moves work to a background thread and then updates Swing components directly from that thread. The freeze disappears, but a data race appears.

Background work should produce data, not mutate Swing. Use \api{SwingUtilities.invokeLater}, a task service, or a Corusco Swing task boundary to return to the EDT. If the screen may close before the result arrives, the task must also handle cancellation or stale-result suppression.

Unsafe background updates are harder to diagnose than visible freezes because
they can appear intermittent. A label update from a worker thread may seem to
work for months. A table model mutation from a pool thread may fail only while
the user is sorting. A validation result may overwrite newer input only under
specific timing. The absence of an exception is not evidence of safety. The
threading contract is violated even when the current run looks fine.

Corusco task and binding helpers are useful because they make the handoff a
named boundary. A task body produces a value. A callback executor returns to the
EDT. A generation counter rejects stale results. A scope cancels or closes
relationships when the screen goes away. These are ordinary engineering
objects, not ceremony. They record the policies that ad hoc threads usually
forget.

The background boundary should also separate computation from presentation.
Returning a list of rows is different from returning a ready-made
\api{JTable}. Returning validation facts is different from setting a red border
on a text field. Returning an image is different from installing it in a label
whose owner may already be disposed. This distinction is not academic. It keeps
worker code testable without Swing and keeps EDT code small enough to audit for
threading safety.

Generation counters are a simple but important habit. The screen increments a
counter when a new request starts. The result carries the counter it was born
with. On completion, the EDT callback compares that generation with the current
generation before publishing state. If they differ, the result is stale. This
pattern applies to search-as-you-type, uniqueness validation, background table
filters, preview rendering, and thumbnail loading. Cancellation is even better
when the work can actually be stopped, but stale-result suppression is still
needed because not every operation can be interrupted immediately.

Structured concurrency, virtual threads, and modern executors make it easier to
write the worker side clearly. They do not remove the need for Swing-specific
handoff. A structured task can still complete after a tab closes. A virtual
thread can still mutate a component illegally. A task scope can still outlive
the screen if it is owned by the wrong object. Good practice is to make the
worker ownership and the Swing ownership meet at one narrow boundary: submit,
observe, cancel or detach, publish on the EDT, close with the screen.

\section{Index-based tables}

A table column view index is a current visual coordinate. It is not a stable application id. Persisting column 3 because it currently displays "Status" breaks when columns move, localization changes, or a hidden column returns.

\begin{coruscobox}{Typed Table Descriptors}
Corusco table descriptors give columns stable keys, resource keys, value types, defaults, and persistence ids. Swing still reorders and sorts, but application meaning no longer depends on current view indexes.
\end{coruscobox}

Bad table code often combines several mistakes: reading selected view row as model row, using localized header text to identify columns, rebuilding the entire model for small changes, and storing column widths by visual order. These mistakes do not always fail in demos. They fail when users sort, filter, localize, hide columns, or keep preferences across versions.

The table mistakes share one cause: using current presentation coordinates as
stable meaning. View row 7 is not a customer. Column 3 is not Status. Header
text is not a persistence id. A selected index is not a selected object. These
shortcuts survive until the table becomes a real table: sortable, filterable,
customizable, localized, persisted, and large enough that users care about
keeping their place.

Good table practice is identity first, coordinates second. A row should have a
business identity when selection or refresh continuity matters. A column should
have a descriptor identity when resources, persistence, validation, or tests
refer to it. View indexes should be converted at the table boundary. Model
indexes should be used only where the model is actually addressed. The code may
be a few lines longer, but the line that converts coordinates is a visible
warning against a whole class of bugs.

Table refresh is another ownership test. If the model calls
\api{fireTableDataChanged} for every small change, the table is told that most
state may be obsolete. Selection, sorter assumptions, editors, and repaint
regions all suffer. If the model fires precise row and cell changes, Swing can
keep more of the user's visual context. Corusco observable lists and table
models encourage that precision because row changes are named before they
become \api{TableModelEvent}s.

A table that stores column widths by current visual order is also borrowing
trouble. Users can move columns. Future versions can add columns. Localized
headers can change. Hidden columns can return. Store widths, visibility, order,
and sorting by stable column ids. Then the table can evolve without treating
user preferences as meaningless integers from an old screen.

The index distinction deserves repetition because it is the source of many
polite bugs. Swing has a model coordinate system and a view coordinate system.
Sorting, filtering, column movement, and hidden columns change the view without
changing the model's application meaning. A selection listener receives a view
row because the user selected something visible. A table model stores data by
model row because the data source has its own order. A table column model
describes visible columns in view order. A descriptor describes application
columns by stable identity. Good table code crosses these coordinate systems
only at named boundaries.

The simplest rule is to convert as late as possible when reading from the UI
and as early as possible when driving the UI from the model. If a user clicks a
view row, convert it to a model row before screen-presenter code asks for the selected
object. If screen-presenter code wants to select an object by id, resolve the id in
the model and then convert the model row to the current view row before touching
the selection model. Code that mixes conversions in the middle of a business
method tends to accumulate off-by-one defects that appear only after a sorter
or filter is installed.

Editing adds another layer. A cell editor edits visible state, but commit must
land in the model fact that owns the value. If the table is sorted while an
editor is active, or if a background refresh replaces the row, the edit policy
must be explicit. Should the editor commit before refresh? Should it cancel?
Should selection follow the row identity? Should validation attach to the
column descriptor and row id rather than to the current view coordinates? These
questions are easier to answer before the table has a production user's saved
layout and ten thousand rows.

Corusco table descriptors are not merely a convenience for generated columns.
They are a way to avoid smuggling application meaning through visual
coordinates. A descriptor can carry the column key, value type, header resource,
tooltip resource, width defaults, persistence id, renderer policy, editor
policy, and validation target. The table remains a Swing \api{JTable}; the
application no longer has to pretend that the current third visible column is a
durable concept.

\textbf{Renderer allocation.}
A renderer is a stamp. It is reused for many cells. It must be fast, complete, and stateless with respect to previous calls. Allocating expensive objects or remembering per-row state inside a renderer creates performance and correctness failures.

The renderer should set all visual properties needed for the current cell: text, icon, foreground, background, font, border, tooltip if appropriate, and opacity. It should not assume a property left over from the previous cell is still correct.

Renderer code should be reviewed as hot-path code. It may run while scrolling,
resizing, selecting, sorting, or repainting after a model update. It should not
open files, query services, create expensive formatters, allocate icons, or
perform permission checks that belong to the model. Precompute or cache stable
presentation facts outside the renderer. Let the renderer stamp the current
cell cheaply and completely.

Completeness is as important as speed. If a renderer sets the foreground red for
error cells, it must set the foreground back for normal cells. If it sets an
icon for warnings, it must clear the icon for non-warnings. If it changes
opacity, font, border, alignment, or tooltip, it must make the ordinary case
explicit too. Reused renderers do not remember "default" for you; they remember
whatever the last cell did.

Use the model to carry row facts that should be tested. A severity value,
validation problem, stale marker, or dirty flag belongs in row state or an
associated presentation model. The renderer can translate that fact into color,
icon, or tooltip. If the renderer computes the fact itself, tests must paint a
table to verify business meaning, and performance becomes unpredictable.

A renderer should also be boring. Boring renderers are good renderers because
they do the same small set of operations for each cell: read already prepared
facts, set component state completely, return the stamp. The interesting
decisions happen before painting. Is the row stale? Is the amount negative? Is
the record locked? Which severity applies? Which icon should represent that
severity? Those decisions can be tested as model or screen-presenter logic. The
renderer then becomes a measured projection of facts to pixels.

Allocation discipline does not mean avoiding every object at all costs. It
means avoiding accidental work in a path called repeatedly during interaction.
A renderer that creates a new date formatter, image icon, border, color palette,
or regular expression matcher for every cell will perform acceptably in a tiny
example and poorly in the screen that matters. Stable resources should be
created once, shared through descriptors or renderer factories, and invalidated
only when their true inputs change: Look-and-Feel, font, scale, theme,
localization, or resource set.

Completeness can be reviewed mechanically. For every property changed in the
exceptional case, find the ordinary case that resets it. For every row-specific
decision, find the model fact that carries it. For every expensive object, find
the cache or descriptor that owns it. For every tooltip, find the policy that
composes it with validation and help. A renderer that passes this review is not
only faster. It is less likely to carry visual residue from the last cell into
the next one.

\begin{pitfallbox}{Renderer Memory}
A table renderer that stores "current row" in a field is almost always wrong. The same renderer instance is reused. Put row-specific state in the model or compute it during the rendering call.
\end{pitfallbox}

\textbf{Focus hacks.}
Focus bugs invite hacks. A developer adds a mouse listener that calls \api{requestFocus} everywhere. Another adds an \api{InputVerifier} that refuses to let focus leave an invalid field. Another consumes Enter in a key listener. The screen begins to feel hostile.

Good focus practice begins with workflow. Which components are focusable? What is the traversal order? Which commands should be available from the keyboard? When validation fails, where should focus go? Which shortcuts belong to text editing and which belong to the window?

Use focus traversal policy when traversal order matters. Use key bindings for commands. Use validation focus only when it helps the user recover. Never trap users in a field merely to prove that validation exists.

Focus hacks often begin after a real user complaint. The tab order is wrong, a
button does not respond to Enter, a custom component does not show focus, or an
invalid field is hard to find. The quick repair is to request focus manually in
several listeners. The durable repair is to decide the focus policy: initial
focus, traversal order, default command, cancel command, field-level validation
display, and recovery after OK fails. Once that policy exists, the code becomes
less frantic.

Keyboard completeness belongs with focus. A screen is not complete merely
because every operation has a button. Can the user reach every field without a
mouse? Can the user invoke primary commands from keyboard shortcuts or default
buttons? Do text components keep their editing shortcuts? Does Escape cancel or
close according to dialog policy? Does F1 open relevant help? These questions
are usability questions and architecture questions at the same time because
they determine where commands and input maps live.

Input verifiers are not a general validation architecture. They can be useful
for local focus-transfer checks, but they are easy to abuse. If an invalid field
prevents focus from leaving, the user may be unable to inspect another field,
copy information, press Help, or cancel gracefully. Prefer validation models,
problem summaries, and command-state rules. Move focus to a problem when it
helps recovery; do not use focus as a cage.

Focus should also be considered with custom components. A custom canvas that
handles keyboard shortcuts must be focusable, paint a visible focus indicator,
advertise accessible information, and avoid stealing traversal keys that belong
to the container. A table cell editor must cooperate with the table's editing
policy. A dialog default button must not consume Enter intended for a
multi-line text component. These details look small, but they decide whether a
screen can be used efficiently without a mouse.

Corusco command metadata helps here because it separates "what operation is
available" from "which widget currently presents it." Swing key bindings still
do the real keyboard work. The command identity, enablement, help topic,
shortcut resource, and disabled reason can come from a shared command model.
Then a toolbar button, menu item, popup item, and root-pane binding agree
without asking each other for state.

\textbf{Tooltips as a shared resource.}
Tooltips are easy to overwrite. Help text wants the tooltip. Validation wants the tooltip. Disabled command reasons want the tooltip. Generated field metadata may want the tooltip. The last writer wins, and all previous meaning disappears.

The better model is composition. Tooltip content can have contributors: static help, validation problems, disabled reasons, and perhaps formatting hints. Corusco's tooltip and behavior concepts exist because tooltip ownership is otherwise invisible.

Tooltip composition should have priority rules. An error may deserve first
position. Static help may follow. A disabled reason may replace command help
when the command is unavailable. A formatting hint may be useful only while the
field is focused. If these priorities are not modeled, whichever listener writes
last becomes the tooltip policy. That is not a policy; it is scheduling.

Tooltips are also not the only presentation of help or validation. A screen
reader may need accessible descriptions. A dialog may need a validation summary.
A status strip may show hover help. A command palette may show disabled
reasons. Once these outputs exist, tooltip text should be generated from shared
facts, not used as the source of truth. Corusco problem sets, resource keys, and
tooltip policies make that sharing practical.

Tooltip policy is also a good example of reader empathy. A beginner can
understand "set the tooltip text" in one line. The intermediate programmer must
then learn that several features want to write the same property. The advanced
programmer should see the property as an output of a policy, not as a storage
place. This is the book's general progression: first show the Swing mechanism,
then show the failure mode under scale, then introduce the Corusco concept that
names the missing owner.

\textbf{Hard-coded resources.}
Hard-coded strings begin as convenience and end as archaeology. A button text appears in a toolbar, menu, popup, and test. A table header string is used as a persistence key. A validation message is copied into a tooltip and a dialog. When localization arrives, the copies disagree.

Good practice is to separate stable identity from visible text. A command has an action key and resource keys. A table column has a column key and header resource. A help topic has a help key. Visible strings are resolved through resources.

Hard-coded resources hurt even before localization. They make tests brittle
because tests search for current English text. They make screenshots harder to
review because the same command may appear with slightly different wording in
two places. They make help inconsistent because a button says "Remove" while a
menu says "Delete". Resource keys are not only about translation; they are about
one screen vocabulary.

Typed resource keys add another guard. A label resource, tooltip resource,
message resource, icon resource, and help topic are not interchangeable merely
because they can all be represented by strings somewhere. The compiler can help
when the application gives these roles distinct types. Generated companions can
then make missing or duplicated resource identities visible during review.

Resource identity should appear before resource lookup in the design. A command
is not named by the English text currently displayed on its button. A help page
is not named by its current title. A table column is not named by a localized
header. Stable identities let user preferences, generated tests, help routing,
accessibility descriptions, and screenshots survive wording changes. Visible
text can improve over time without becoming a schema migration.

The strongest resource practice is to make missing vocabulary visible early.
Generated companions can list expected keys. Tests can assert that command
resources exist. A screenshot can reveal awkward text, but it should not be the
first place a missing resource is discovered. When resources are treated as a
late decoration, the code tends to hard-code "temporary" strings that become
part of tests and user expectations. When resources are part of command and
descriptor identity, wording remains editable.

\textbf{Preferred-size abuse.}
Setting preferred sizes everywhere is a layout smell. Some components have legitimate fixed or preferred dimensions: icons, swatches, small numeric fields, fixed preview thumbnails. But using preferred sizes to force every field and panel into a screenshot shape fights layout managers, Look-and-Feel changes, fonts, localization, and high-DPI scaling.

Prefer layout constraints. Let the label column take label width. Let the field column grow. Give the main table row vertical growth. Put large custom views in scroll panes with correct preferred sizes. Use component preferred size to describe the component's natural content, not to micromanage the parent.

The component chapter made the positive rule explicit: a component's preferred
size should describe its natural content; the parent layout describes the
relationship among components. Preferred-size abuse flips those roles. It asks
children to compensate for a parent layout that has not been designed. The
result may look acceptable until a font, theme, locale, scale, or window size
changes. Then every hard-coded preferred size becomes a small refusal to adapt.

There are legitimate preferred sizes. A custom canvas may report a virtual
world. A color swatch may be fixed. A scroll pane may have a preferred viewport
size. The smell is not "preferred size exists"; the smell is "preferred size is
being used to force layout because the parent constraints are unclear." During
review, ask whether the size belongs to the component's nature or to one
screen's geometry. If it is screen geometry, put it in layout policy.

This distinction matters especially with MigLayout. MigLayout gives the parent
enough vocabulary to express relationships directly: columns, rows, growth,
alignment, wrapping, related gaps, unrelated gaps, and docking. A developer who
sets preferred sizes on children to obtain those relationships is bypassing the
layout tool. The result may pass one screenshot and fail the moment labels are
longer, fonts are larger, or the user stretches the window. The better practice
is to let components tell the truth about content and let MigLayout tell the
truth about screen geometry.

Scroll panes are the most common place where this rule becomes visible. The
view's preferred size describes the virtual world or natural content. The
viewport's extent describes what is currently visible. The parent layout
decides how large the scroll pane should be in the surrounding screen. Calling
\api{setPreferredSize} on the scroll pane to force a table height may be
reasonable for a compact dialog, but it should be a deliberate viewport policy,
not a hidden attempt to make rows fit a screenshot. The Components and MigLayout
chapters provide the mechanics; good practice is choosing the owner.

\section{Global state and client properties}

Swing client properties are useful extension points. They are also an easy place to hide untyped state. A Look-and-Feel client property that controls a documented visual option is fine. A random object graph stored under a string key because it was convenient is not.

Global singletons have a similar risk. A global help service, resource service, or diagnostic registry may be appropriate. A global "current customer panel" is not. Prefer passing services through screen presenters, behavior contexts, or application shells.

The question for any global is whether it represents application infrastructure
or current screen state. A resource bundle service is infrastructure. A help
router may be infrastructure. A task executor owned by the application shell may
be infrastructure. The currently selected customer, current dialog, current
table row, and current validation problem are not infrastructure. They are
screen state, and making them global creates invisible coupling between
screens.

Client properties deserve the same discipline on a smaller scale. Look-and-Feel
libraries often document client properties that alter component appearance or
behavior. Those are reasonable when used deliberately. Application code can also
use client properties to attach component keys or small metadata. The bad habit
is using string-keyed properties as a hidden object graph. If a property becomes
essential to workflow, give it a typed owner or descriptor.

Global state also makes tests lie. A test that passes because a previous test
left a global current user, resource bundle, or selected customer is not a
screen test; it is an order-dependent accident. Passing services explicitly or
through a narrow application context keeps tests honest. Corusco typed keys and
descriptors can reduce the temptation to use global registries for every fact.

Client properties can still be a disciplined boundary when the key is stable
and the value is small. Component keys, help topics, test identifiers, and
Look-and-Feel hints are examples. The danger begins when a client property
becomes a private service locator or a side channel between unrelated code.
String keys make such channels hard to search and easy to collide. If a
property must be used, give the key a constant, document the expected value
type, and keep ownership close to the component or behavior that consumes it.

Global registries need a similar test. A registry that maps command ids to
command descriptors may be application infrastructure. A registry that records
the current selected row because "many places need it" is usually a hidden
model. The first kind of global reduces duplication. The second kind makes
screens order-dependent. The difference is whether the global describes stable
application vocabulary or mutable screen state.

\textbf{Dialogs with copied button logic.}
Many Swing applications have dozens of dialogs. Each dialog reimplements OK, Cancel, Escape, Enter, dirty confirmation, validation focus, and active editor commit. The copies drift. One dialog loses edits because the active table cell editor did not commit. Another closes on Escape without dirty confirmation. Another validates but does not focus the failing field.

Good practice is to standardize dialog workflow. Corusco's dialog helpers exist for this reason. Layout remains local; lifecycle and commit/cancel policy should be consistent.

Copied dialog logic is expensive because users experience dialogs as a family.
If Enter commits one dialog and inserts a newline in another, the difference
should be intentional. If Escape cancels one dirty dialog without confirmation
but protects another, the difference should be explainable. If one dialog
commits active table editors and another loses the final edit, users lose trust.
The workflow is common enough to deserve common code.

The dialog body may vary greatly. One dialog may contain a two-column MigLayout
form. Another may contain a table and detail pane. Another may contain a
scrollable validation review. But the shell policy remains familiar: active
editor commit, validation, focus or summary of blocking problems, OK result,
Apply baseline, dirty cancel confirmation, Escape, default button, owner
window, disposal, and cleanup. Standardizing the shell lets the body remain
specific without copying fragile ceremony.

A dialog should also be testable without becoming modal theatre. Form model
tests prove validation and result creation. Dialog controller tests prove OK,
Cancel, Apply, dirty confirmation, and active-editor commit. A small EDT test
proves the Swing bindings and keyboard actions. A screenshot proves the visual
states. If every dialog rule is tested only by clicking a modal window, the
suite becomes slow and imprecise.

Dialog consistency is not sameness of layout. It is sameness of promises.
Users should be able to predict which key commits, which key cancels, whether
dirty changes are protected, whether validation blocks closing, where the first
problem is shown, and whether a background operation can still complete into a
closed dialog. The body can be a simple form, a table editor, a wizard page, or
a rich custom component. The shell policy should remain recognizable.

This is a strong use case for Corusco because dialog workflow touches many
feature sets at once. Form models own parse and validation state. Commands own
OK, Apply, Reset, and Cancel availability. Problems own severity and targets.
Bindings own field synchronization. Task services own asynchronous validation
or save. Scopes own cleanup. Generation can remove repeated wiring once the
dialog facts have names. The point is not that every dialog must be generated.
The point is that every nontrivial dialog should stop inventing the same
workflow in private.

\textbf{Leaking screen lifetimes.}
A screen closes visually but remains alive because a timer, service listener, property-change listener, or task callback references it. This is one of the most common long-session Swing failures. It is also one of the hardest to diagnose after the fact.

Make cleanup visible. A screen activation creates a scope. Bindings, subscriptions, table controllers, timers, and task handles are added to it. Closing the screen closes the scope. Tests should verify at least the common cleanup path.

Leaks are often behavior bugs before they are memory bugs. A reopened screen
receives duplicate events. A closed dialog updates a status line. A timer keeps
repainting a component that is no longer visible. A background task writes rows
into a new screen that happens to reuse the same model. These symptoms should
trigger a lifecycle review immediately. Do not wait for a heap dump to prove
that the old object graph is still reachable.

The cleanup owner should be created at the same boundary as the relationships.
If a panel installs bindings when it is constructed, the panel or its controller
should close them. If a tab installs subscriptions when activated, deactivation
or close should release them according to policy. If a dialog starts tasks, the
dialog controller should cancel or detach them when the dialog closes. A single
application shutdown hook is too late for screen-local relationships.

Weak listeners are not a general cleanup strategy. They can be useful when a
long-lived source must not retain an observer, but they do not express workflow
ownership. A weak listener can vanish too early, remain because another
reference exists, or hide the fact that a relationship should have been closed.
Prefer deterministic scopes when the lifecycle boundary is known.

Lifetime review should include the sources that live longer than the screen:
application event buses, services, table models shared across views, timers,
executor callbacks, static caches, global diagnostics, document models, and
Look-and-Feel listeners. A listener from a short-lived screen to a long-lived
source is a retention path unless it is removed. A task callback that captures
\api{this} is a retention path until it completes or is canceled. A timer that
fires into a disposed component is a behavior defect even if memory pressure is
not yet visible.

There is a useful symmetry between installation and cleanup. Code that installs
a relationship should usually record its cleanup in the same small block. If
the installation is generated, the generated companion should also own the
disposal hook. If a behavior attaches listeners, it should return a binding or
handle that can be closed. If a task is started by a screen, the screen's scope
should decide whether close cancels it, detaches from it, or lets it complete
into a shared model. The cleanup policy is part of behavior, not garbage
collection trivia.

Tests can catch lifetime mistakes cheaply. Create the screen or binding, attach
it to a model, close its scope, mutate the model, and assert that the old view
does not change. For tasks, start a controlled task, close the screen, complete
the task, and assert that no disposed component is touched. These tests are
small because the lifetime boundary is explicit. Without a scope or handle, the
test has to reverse-engineer private listeners and becomes too awkward to write.

\textbf{When plain Swing is good practice.}
Good practice does not mean using Corusco everywhere. Plain Swing is good practice when the behavior is local, obvious, and has no durable identity. A one-off label update in a small diagnostic panel may not need a model. A tiny utility window may not need generated descriptors. A local renderer for a special table column may be clearer than a framework abstraction.

The decision point is not "framework or no framework". It is "does this fact need a stable owner?" If the answer is yes, name it. If the answer is no, keep the plain code plain.

Plain Swing is also good practice for learning. Before a binding abstraction is
useful, the developer should understand the document listener it replaces.
Before a table descriptor is useful, the developer should understand model and
view indexes. Before a busy overlay binding is useful, the developer should
understand the EDT and repaint. This book keeps plain Swing baselines because
framework use without baseline knowledge becomes superstition.

The boundary can move over time. A local diagnostic panel may start plain and
remain plain. A utility dialog may begin with direct listeners and later gain a
form model when validation grows. A table may begin with a small
\api{AbstractTableModel} and later adopt descriptors when persistence,
resources, editing, and validation appear. Migration is not failure; it is a
sign that the screen's responsibilities became clearer.

\textbf{Corusco as a discipline amplifier.}
Corusco cannot rescue careless Swing code by being present on the classpath. It helps when the code accepts its discipline: typed keys for stable identity, models for screen state, commands for operations, descriptors for metadata, problems for validation, bindings for synchronization, behaviors for reusable component capabilities, and scopes for cleanup.

This is why migration should start with review. Find the hidden facts first. Then choose the Corusco feature that gives each fact an owner.

Corusco works best when the team agrees on the discipline it represents. Values
represent changing facts. Commands represent operations. Problems represent
validation and parse issues. Descriptors represent stable metadata. Bindings
represent synchronization. Behaviors represent reusable component capabilities.
Scopes represent lifetimes. Generation removes repetition once the source facts
are stable. This vocabulary is small enough to use in code review and precise
enough to prevent most ad hoc listener growth.

A poor migration starts by replacing syntax. A good migration starts by naming
facts. For a problematic form, name the fields, dirty state, touched state,
problems, command state, and lifecycle scope. For a problematic table, name row
identity, column identity, row source, selection, sorter policy, renderer facts,
and table state. For a problematic background operation, name busy state, task
handle, cancellation, generation, success, failure, and stale-result policy.
Only then choose the Corusco APIs that fit.

This approach also protects the codebase from over-abstraction. Some hidden
facts become plain model fields. Some become Corusco values. Some become
commands. Some become tests. The point is not to maximize library surface area.
The point is to make screen meaning reviewable. Corusco amplifies discipline
that already exists; it cannot replace the decision to be disciplined.

The most effective migration sequence is therefore modest. First name the
facts without changing the UI. Second move repeated facts out of widgets and
into a screen presenter or model. Third introduce Corusco values, commands, problems,
descriptors, and scopes where they match those named facts. Fourth replace
repeated Swing wiring with bindings or behaviors. Fifth add generation only
after the source facts are stable enough that generated code will remove
duplication rather than hide uncertainty. This sequence respects existing code
and still moves it toward a disciplined architecture.

It is also acceptable to stop partway. A screen may gain command descriptors
and resource keys while keeping a hand-written form. A table may gain column
descriptors while leaving a specialized renderer manual. A dialog may gain a
standard OK/Cancel controller while keeping local validation. The review
question is whether the remaining manual code is still local and obvious. If it
is, keeping it plain is a good practice. If it is copied, stateful, hard to
test, or shared by several presentations, it probably wants a stronger owner.

\section{Case studies and drills}

\textbf{The save button that lies.}
The Save button remains enabled while a background save is running. The user clicks twice. The second request races the first. The fix is not merely disabling the button in the listener; it is making saving state part of the command model.

The good version has a command whose enabled state is derived from dirty,
committable, permission, and busy facts. When saving begins, busy state changes
and every presentation of Save changes together: toolbar button, menu item,
shortcut, default button, and tooltip. If saving fails, the command state and
problem or status model explain recovery. If saving succeeds, the dirty
baseline is accepted. The button no longer has to lie because it no longer owns
the truth.

\textbf{The status label as event log.}
A status label is used for validation, progress, hover help, and errors. Messages overwrite each other unpredictably. The fix is to decide which status source has priority, or to model separate status regions.

A single status line can still be useful, but it needs a policy. Validation
errors may deserve a problem summary. Hover help may be transient and low
priority. Background failure may deserve a persistent message until dismissed.
Progress may deserve its own indicator. Once priorities are named, the UI can
compose a status presentation rather than letting listeners race to write text.
The same facts can also feed tooltips, accessibility descriptions, and tests.

\textbf{The table that forgets itself.}
Every refresh calls \api{fireTableDataChanged}. Selection clears, column widths drift, sorter state resets, and repaint cost spikes. The fix is precise list/table events and stable row identity.

The better version asks what kind of refresh occurred. Were rows inserted,
removed, replaced, or reordered? Did a row keep the same identity but receive a
new value? Should selection follow identity, current view position, or nothing?
Should column state survive because column ids are stable? A table refresh is a
small reconciliation problem, not a button that says "everything changed".
Corusco observable lists and descriptors give the reconciliation named inputs.

\textbf{The dialog that cannot be tested.}
Validation is implemented entirely by labels in a modal dialog. The test must click through the UI. The fix is a form model with validation tests, plus a small dialog binding test.

The model-backed dialog has a different proof structure. A form model test
checks parse state, problems, dirty state, touched state, and result creation.
A command or dialog-controller test checks OK and Cancel policy. An EDT test
checks that fields and labels are bound to the right model facts. A screenshot
checks layout and visible validation states. Each test fails near the defect.
The modal window is no longer the only way to know whether the dialog works.

\textbf{The renderer that worked with ten rows.}
A renderer formats dates by creating a formatter and loading an icon on every cell paint. It passes with ten rows and stutters with ten thousand. The fix is cached immutable resources, cheap rendering, or a measured optimized renderer.

The review should ask what work happens per visible cell and what work happens
per row update. Date formatters, icons, severity decisions, and permission flags
can often be prepared outside the renderer. The renderer then sets text, icon,
colors, border, and tooltip from already available facts. Measure with
realistic row counts, because renderer defects hide in small demos and appear
as "Swing is slow" in production.

\textbf{The component that fights layout.}
A custom component calls \api{setPreferredSize} from outside whenever data
changes. Sometimes the scroll bars update, sometimes they do not. The fix is to
move geometry into the component's \api{getPreferredSize}, call
\api{revalidate} when the preferred size changes, and let the parent layout use
that truthful size. This case connects the component chapter to the practice
chapter: bad layout practice is often a component contract violation disguised
as a parent problem.

\textbf{The background validation race.}
A uniqueness check starts on every edit. The user types "Ada", then "Adal".
The slower result for "Ada" returns last and marks the current value invalid.
The fix is not merely "use a background thread"; the screen also needs
generation or cancellation policy. A task result must prove it still belongs to
the current field value before it publishes problems. Otherwise the code is
fast and still wrong.

\textbf{The help shortcut that opens the wrong page.}
F1 always opens the application's front page because no component has a stable
help topic. Users learn not to use help. The fix is to give fields, commands,
tables, and screens typed help topics and route F1 from current focus or command
context. Swing installs the key binding; Corusco metadata supplies the target.
The result is not more decorative help. It is a screen whose meaning can be
found from the keyboard.

\textbf{The leak that looked like validation.}
After a dialog is opened three times, each invalid edit shows three identical
messages. The validation rule is not wrong. Three old listeners are still
subscribed. The fix is lifecycle ownership: the dialog activation creates a
scope, every binding and subscription enters it, and closing the dialog closes
the scope. A cleanup test then mutates the model after close and proves that the
old view no longer reacts.

\textbf{Review checklist.}
\begin{itemize}
\item A listener translates events; it does not secretly own multiple screen facts.
\item EDT work has bounded, predictable latency.
\item Background work returns to Swing through an explicit EDT boundary.
\item Table meaning uses descriptors or keys, not localized headers or view indexes.
\item Renderers are fast, complete, and reusable.
\item Focus policy is deliberate and keyboard-friendly.
\item Tooltips compose help, validation, and disabled reasons instead of overwriting blindly.
\item User-facing strings have resource identities when they outlive a demo.
\item Preferred sizes describe component content; layout constraints describe parent geometry.
\item Every installed listener, subscription, timer, task, binding, and controller has a cleanup owner.
\end{itemize}

The most productive way to use this chapter is not to memorize every smell. It
is to practice translating symptoms into ownership questions. When a button is
wrongly enabled, ask who owns command state. When validation messages disagree,
ask who owns problems. When a table forgets selection, ask who owns row
identity. When a screen leaks, ask who owns installed relationships. These
questions turn vague unease into a review path.

Good Swing code is often ordinary-looking code with good boundaries. It still
creates panels, buttons, fields, tables, actions, and renderers. It still has
listeners. It still uses layout managers. The difference is that component
callbacks are narrow and state lives somewhere named. A reader can find the
model, find the command, find the validation rule, and find cleanup without
reverse-engineering a chain of side effects.

Bad practice often survives because it is useful at first. An inline listener is
the fastest way to prove a button works. A hard-coded string is the fastest way
to see text on screen. A full table reset is the fastest way to refresh rows.
The transition from prototype to application is the moment to retire those
shortcuts. Keeping them after the screen has durable workflow is what turns
convenience into maintenance debt.

Do not replace one kind of obscurity with another. A framework abstraction that
hides every listener but gives no clear model is not an improvement. Corusco's
benefit is strongest when it names facts that already exist: field keys,
commands, problem codes, descriptors, bindings, scopes, and task handles. If a
screen has not identified those facts, adding library calls will only move the
confusion.

The review order is therefore deliberate. Start with user-visible workflow.
Which operations can the user invoke? Which state enables them? Which errors
block them? Which data can refresh? Which selection should survive? Then inspect
the Swing mechanics. Which components present those facts? Which listeners
translate events? Which bindings synchronize state? Which scopes close them?
This order prevents code review from getting lost in widget details before the
screen's promises are clear.

Testing follows the same order. Model tests prove validation, dirty state,
command enablement, and row identity. EDT tests prove that components are bound
to the right facts and clean up correctly. Screenshot tests prove that visual
states are present and stable enough for the book. Manual testing remains
useful, but it should not be the only proof that the screen obeys its own
contracts.

Finally, remember that a good practice can become bad when applied without
scale. A tiny diagnostic popup may not need a generated form model. A one-cell
renderer may not need a descriptor system. A local listener may be clearer than
a named command for a one-off component. The rule is not "use more framework".
The rule is "use enough structure that ownership is visible at the scale of the
screen."

Corusco should make disciplined Swing cheaper, not make every program heavier.
When it removes repeated wiring, stabilizes identity, or makes cleanup explicit,
use it confidently. When plain Swing already expresses a local fact perfectly,
leave the code plain. The mature style is selective: explicit where meaning
outlives a component, simple where behavior is genuinely local.

The examples in this book should be read with the same selectivity. A compact
example may use a direct listener because the chapter is teaching painting or
layout, not command architecture. A later example may replace that listener
with a command because the operation now has several presentations. This is not
inconsistency. It is the progression from local demonstration to application
structure, and it is the reason the book keeps both plain Swing baselines and
Corusco treatments.

When a chapter feels like a list of prohibitions, return to the positive rule:
make the screen's promises explicit. The user should know what can be done, why
it cannot be done, what data is being shown, what is still loading, and what
will happen when the screen closes. The programmer should be able to find the
same answers in named models, commands, descriptors, tasks, and scopes.

A good practice review can be written as a small story. The screen opens with
known state. Components show that state. The user changes one fact. The model
records the change. Derived values update command enablement, validation, and
status. Slow work leaves the EDT and returns through a documented boundary.
Tables preserve identity. Renderers stamp visible cells cheaply. Tooltips and
help compose from shared metadata. Closing the screen closes installed
relationships. If the story can be told without mentioning "and then this
listener also...", the design is probably healthy.

The same story should be visible in code. A screen presenter or model names the facts.
A view constructs components and layout. Bindings or behaviors synchronize
component state with model state. Commands own operations. Descriptors own
stable metadata. Scopes own lifetimes. Tests exercise the layers in the same
order. This structure is not ornamental; it is how a Swing screen remains
readable after the third feature request.

There is also an aesthetic reward. Disciplined Swing code feels calmer. The
view is no longer a tangle of urgent callbacks. The model no longer depends on
which button was clicked. The table no longer forgets itself. The dialog no
longer invents OK and Cancel policy. The code still uses ordinary Swing, but it
has fewer secrets. That is the practical meaning of "good practice" in this
book.

Use this chapter as a diagnostic companion to the rest of the book. When a later
chapter introduces values, lists, forms, commands, tasks, bindings, tables, or
generation, ask which bad habit the feature is meant to make unnecessary. The
answer should be concrete. A Corusco feature earns its place when it removes a
real source of drift, duplication, latency, hidden state, or lifecycle risk.

\section{A practice review in the field}

The most useful practice review is neither a style debate nor a rewrite plan.
It is a short investigation of one screen that already matters. Choose a screen
with enough behavior to be interesting: a search panel, an editor dialog, a
table with actions, a background refresh view, or a form with validation. Start
from the user's view of the screen. Write down the operations, the visible
states, the data that may change, the commands that may be disabled, the
validation problems that may block progress, and the lifetime of the screen.
Only then read the code. This order prevents the review from becoming a tour of
listeners before anyone has stated what the listeners are supposed to preserve.

The first reading should look for owners rather than defects. Where is the
selected object stored? Where is dirty state stored? Where is command
enablement computed? Where are validation problems produced? Where is busy
state represented? Where are resource keys defined? Where are table columns
identified? Where are installed listeners closed? If a fact has a clear owner,
the review can move on. If a fact appears only as a component property, a
listener-local variable, or a repeated expression, mark it. Repeated ownership
questions usually reveal the real design problem faster than a line-by-line
search for mistakes.

The second reading should follow one user story from input to visible result.
For example: the user edits a field, the field becomes dirty, validation runs,
Save becomes enabled or disabled, the tooltip changes, the status line changes,
and OK either commits or focuses the first problem. In a healthy screen, this
story can be told with named facts. In an unhealthy screen, the story becomes a
tour of incidental callbacks: the document listener sets the label, the focus
listener sets the tooltip, the button listener checks a flag, the status label
is overwritten by a timer, and the dialog closes through a copied action. The
code may still work today, but the story shows why it will be difficult to
change.

The third reading should follow the slow path. What happens if the data source
is slow, the server rejects a save, the user edits again while validation is
running, or the screen closes before a worker completes? Many Swing defects are
not algorithmic. They are missing answers to timing questions. A screen without
busy state invites duplicate commands. A screen without stale-result handling
publishes old validation. A screen without cancellation keeps work alive after
the user has moved on. A screen without failure state treats exceptions as log
events rather than user-visible facts. The repair is not merely "use a thread";
it is to make latency part of the model.

The fourth reading should examine identity. Look for view row indexes that
escape table code, column numbers stored in preferences, localized header text
used as keys, selection represented only by current row, resource text copied
into tests, and help topics derived from widget labels. These mistakes often
look harmless because the screen has one locale, one table order, and one demo
data set. They become expensive when the user sorts, filters, moves columns,
customizes state, reopens the application, or changes language. Stable identity
is not a large-framework concern. It is the only way to let visible
presentation change without losing application meaning.

The fifth reading should inspect hot paths and visual correctness. Renderers,
painting methods, layout invalidation, tooltips, and table refresh events all
run more often than developers tend to imagine. A renderer that allocates
formatters may be invisible in a test and obvious during scrolling. A custom
component that calls \api{repaint} too broadly may work until it is placed in a
split pane. A table model that fires full resets may pass functional tests and
still destroy user context. A tooltip writer that ignores validation may make a
screen look polished while hiding the reason a command is disabled. Good
practice includes performance because performance is part of the user's
experience of correctness.

After the readings, do not propose a heroic rewrite. Classify findings by the
owner they need. Some findings need plain refactoring: a screen-presenter method, a
small model class, a cached renderer resource, or a cleanup handle. Some need
Corusco concepts: a value, command, problem set, descriptor, behavior, task
service, or scope. Some need generated companions because the repeated metadata
has become stable. Some need tests before code moves. The important discipline
is to match the remedy to the missing owner. Replacing every listener with a
framework call is as unfocused as leaving every fact hidden in widgets.

A good first repair is small and observable. Move one duplicated command
enablement rule into a command state. Give one table column a stable identity
and persist state by that identity. Replace one tooltip overwrite with a
composition policy. Add one scope that owns listener cleanup. Convert one
blocking refresh to a task that publishes on the EDT and ignores stale results.
Each repair should leave a test or screenshot trail. The chapter's advice
becomes credible in a codebase when every step makes the next defect easier to
find, not merely when the code looks more modern.

Good repairs also have a teaching value. A codebase full of Swing screens is
usually maintained by people with different levels of Swing memory. One
developer remembers renderer reuse but not table coordinate conversion. Another
remembers the EDT but not dialog editor commit. Another knows MigLayout but not
client-property discipline. A repair that names the owner and leaves a small
test becomes onboarding material. The next maintainer learns the rule from the
code rather than from a private conversation.

This is one reason the companion examples should remain compact. A book example
does not need every feature of the production screen. It needs the smallest
screen that makes the ownership visible: one bad listener pile beside one
model-backed flow, one table that preserves selection by identity, one renderer
that caches resources, one dialog that commits before validation, one task that
ignores stale completion. Screenshots then become evidence for the prose. They
show that the rule is not only conceptual; it produces visible, repeatable
Swing behavior under FlatLaf and MigLayout.

Reviews should also preserve local clarity. If a small panel has one button
that closes an in-memory preview, an inline listener may be the clearest code.
If a custom renderer for a single visual column is short, complete, and tested,
it need not become a generated artifact. If a dialog has no dirty state, no
background work, and no shared command identity, a simple hand-written OK
action may be good practice. A practice review should not punish simplicity.
It should punish hidden durable meaning. That distinction keeps Corusco from
becoming ritual and keeps plain Swing from becoming an excuse for drift.

The review output should be written in vocabulary the team can reuse. Instead
of "the listeners are messy", write "dirty state is owned by three listeners and
the Save button." Instead of "the table code is fragile", write "row identity is
lost at refresh because selection is stored by view index." Instead of "the
dialog is inconsistent", write "active editor commit and dirty cancel policy are
copied and differ from the standard shell." Such notes are easier to turn into
work items, tests, and chapter examples. They also teach the team how to see
the next screen.

The final review question is whether the screen has a story that survives the
next feature. If validation gains an asynchronous rule, where will stale
results be suppressed? If a menu item is added for Save, where will enablement
come from? If the table gains hidden columns, where will persistence store
state? If help gains per-field topics, where will focus routing find them? If
the screen is reopened many times in one session, what closes the old
relationships? A design that can answer these questions does not need to be
large. It needs to have the right owners.

This field method also explains why the rest of the book alternates between
plain Swing and Corusco. Plain Swing shows the mechanism and keeps the reader
honest about what the platform actually does. Corusco names the owner when the
mechanism becomes repeated, shared, or lifecycle-sensitive. The distinction is
essential. A reader who sees only plain Swing may underestimate the cost of
scale. A reader who sees only Corusco may treat library vocabulary as magic.
Good practice is the bridge: understand the mechanism, identify the boundary,
then choose the smallest owner that keeps the screen truthful.

The same bridge should guide code comments and documentation. Do not comment
that a listener "updates the button"; the code already says that. Comment when
the listener is the boundary between a Swing event and a named model fact, when
a task result is deliberately ignored because its generation is stale, when a
renderer cache is invalidated by Look-and-Feel change, or when a dialog keeps
focus movable despite validation so the user can still reach Help and Cancel.
Such comments record policy rather than restating mechanics. They are rare, but
they are valuable because they tell the next maintainer which practice rule is
being protected.

The opposite is also true: if a practice rule cannot be explained near the
code, it may not be the right rule for that screen. A tiny panel should not
carry a large lifecycle framework only because another chapter praised scopes.
A local calculation should not become an asynchronous task only because the EDT
chapter warned about blocking. The review must keep scale in view. Good
practice is precise, not maximal. It adds ownership where ownership is missing
and leaves local code local where the screen has no durable promise to protect.
That restraint is part of the discipline, and it is what keeps good practice
from turning into superstition.
The reader should leave with judgement, not recipes.
That distinction matters.

\section{Exercises}

\begin{exercisebox}
\begin{enumerate}
\item Find three listeners in an existing screen and write down who owns their removal.
\item Replace one column-index constant with a typed column identity.
\item Profile a renderer by counting allocations per visible repaint.
\item Convert a blocking refresh listener into a background task with EDT completion.
\item Replace duplicated Save enablement logic with one command state.
\item Audit one screen for tooltip overwrites. Decide which feature should compose tooltip content.
\item Pick a dialog and write its OK/Cancel/dirty/validation workflow as a state table.
\end{enumerate}
\end{exercisebox}
